\documentclass[handout]{ximera}
\usepgfplotslibrary{fillbetween}
\input{../preamble}

\outcome{Apply definite integrals}

\title{4.8 Applications of Definite Integrals}


\begin{document}

\begin{abstract}
In this section we use definite integrals to study rectilinear motion.
\end{abstract}

\maketitle



\section{Definite Integrals and Rectilinear Motion}






\subsection{Rectilinear Motion}

Suppose that an object is moving along a straight path.
The position, velocity, and acceleration functions governing the motion of the object 
can be derived from each other in the following ways.\\
If we start with the position function,$s(t)$,
we can compute velocity and acceleration by differentiating:
\[v(t) = s'(t)\quad \text{and} \quad a(t) = v'(t)\]
On the other hand, if we start with the acceleration function, $a(t)$,
we can compute velocity and position by integrating:
\[v(t) = \int a(t) \ dt \quad \text{and} \quad s(t) = \int v(t) \ dt\]
These indefinite integrals will contain a constant of integration, $C$, 
which can be determined by using the initial conditions of the problem.\\

In the following proposition, we see that definite integrals that can be used to compute the displacement 
and distance traveled by an object moving along a straight path.\\

\begin{proposition}
Suppose that an object moving along a straight path has instantaneous velocity $v(t)$, then
\begin{enumerate}
\item the displacement of the object over the interval $[a,b]$ is given by 
\[s(b) - s(a) =\int_a^b v(t) \ dt\]
and
\item the distance traveled by the object over the interval $[a,b]$ is given by 
\[\int_a^b |v(t)| \ dt.\]
\end{enumerate}
\end{proposition}



\begin{example}[example 1]
An object is launched vertically from a height with an initial velocity of $48 ft/sec$.  
\begin{enumerate}
\item Find the displacement of the object from time $t = 0$ to time $t = 2$.
\item Find the total distance traveled from time $t = 0$ to time $t = 2$.
\end{enumerate}

Part (a): We fist find the velocity function:
\[v(t) = \int a(t) \ dt = \int -32 \ dt = -32t + C\]
To find the value of the constant $C$, we use the fact that the initial velocity was given as $48 ft/sec$:
\[v(0) = -32(0) + C = 48,\]
which implies that $C = 48$.  Thus,
\[v(t) = -32t + 48 or 48 - 32t.\]
The displacement from $t= 0$ to $t = 2$, is given by
\begin{align*}
s(2) - s(0) & = \int_0^2 v(t) \ dt \\
 & = \int_0^2 (48 - 32t) \ dt \\
 & = (48t - 16t^2)\bigg|_0^2 \\
 & = 32 \text{ feet}
\end{align*}

Part (b): To find the total distance traveled from $t = 0$ to time $t = 2$,
we need to integrate $|v(t)|$. To remove the absolute value bars, we analyze the 
sign of the velocity function, $v(t) = 48-32t$. Setting $v(t) = 0$, we find that $t = 3/2$ seconds.
For $t < 3/2$, we have $v(t) > 0$ and for $t > 3/2$, we have $v(t) < 0$.\\
Thus,
\[|v(t)| = \left\{
     \begin{array}{rc}
       v(t) & \ \text{if} \  t < 3/2 \\
			 -v(t) & \ \text{if} \ t > 3/2
     \end{array}
   \right.
\]
The distance traveled from $t= 0$ to $t = 2$, is given by
\begin{align*}
  \int_0^2 |v(t)| \ dt &= \int_0^{3/2} |v(t)| \ dt + \int_{3/2}^2 |v(t)| \ dt\\
  &= \int_0^{3/2} v(t) \ dt - \int_{3/2}^2 v(t) \ dt
\end{align*}

We will compute these integrals separately. The first integral is
\begin{align*}
\int_0^{3/2} v(t) \ dt & = \int_0^{3/2}(48 - 32t) \ dt \\
 & = (48t - 16t^2)\bigg|_0^{3/2} \\
 & = 36 \text{ feet}
\end{align*}
The second integral is
\begin{align*}
\int_{3/2}^2 v(t) \ dt & = \int_{3/2}^2(48 - 32t) \ dt \\
 & = (48t - 16t^2)\bigg|_{3/2}^2 \\
 & = 32 - 36 = -4 \text{ feet}
\end{align*}
Subtracting gives the total distance traveled:
\begin{align*}
\int_0^2 |v(t)| \ dt & = \int_0^{3/2} v(t) \ dt - \int_{3/2}^2 v(t) \ dt \\
 & = 36 - (-4) = 40 \text{ feet}
\end{align*}

\end{example}



\begin{problem}(problem 1a)
An object is launched vertically from a height of $24 ft$ with an initial velocity of $64 ft/sec$.  
\begin{enumerate}
\item Find a formula, $s(t)$, for the height of the object at time $t$
 seconds.
\item Find the maximum height reached by the object.
\item Find the displacement of the object from time $t = 0$ to time $t = 3$.
\item Find the total distance traveled from time $t = 0$ to time $t = 3$.
\end{enumerate}
Assume $g = -32 \ ft/sec^2$.\\
$s(t) = \answer{-16t^2 + 64t + 24},$\\
$\text{maximum height} =\answer{88} \ ft,$\\
$\text{displacement} =\answer{48} \ ft,$\\
and
$\text{distance traveled} = \answer{80} \ ft.$
\end{problem}



\begin{problem}(problem 1b)
An object is launched vertically from a height of $5 ft$ with an initial velocity of $16 ft/sec$.  
\begin{enumerate}
\item Find a formula, $s(t)$, for the height of the object at time $t$
 seconds.
\item Find the maximum height reached by the object.
\end{enumerate}
Assume $g = -32 \ ft/sec^2$.\\
$s(t) = \answer{-16t^2 + 16t + 5},$\\
and
$\text{maximum height} =\answer{9} \ ft.$\\


\end{problem}


\subsection{Average Value}
\begin{definition}[Average Value]
If $f(x)$ is continuous on the interval $[a,b]$, then the average value 
of $f(x)$ on $[a,b]$ is given by
\[f_{ave} = \frac{1}{b-a}\int_a^b f(x) \ dx.\]
\end{definition}



\begin{example}[example 2]
Compute the average value of $f(x) = x^2$ on the interval $[0,2]$.\\
The average value is given by 
\begin{align*}
f_{ave} &= \frac{1}{b-a}\int_a^b f(x) \ dx \\
 & = \frac12\int_0^2 x^2 \ dx \\
 & = \frac{x^3}{6} \Bigg|_0^2 \\
  & = \tfrac43
\end{align*}
\end{example}



\begin{problem}(problem 2)
Compute the average value of $f(x) = x(x+1) +1$ on the interval $[0,3]$.
\begin{hint}
Multiply first, then integrate
\end{hint}
\[f_{ave} = \answer{11/2}.\]
\end{problem} 



\begin{example}[example 3]
Compute the average value of $f(x) = \sin(x)$ on the interval $[0,\pi]$.\\
The average value is given by 
\begin{align*}
f_{ave} &= \frac{1}{b-a}\int_a^b f(x) \ dx \\
 & = \frac{1}{\pi}\int_0^\pi \sin(x) \ dx \\
 & = -\frac{1}{\pi} \cos(x) \Bigg|_0^\pi \\
  & = \frac{2}{\pi}
\end{align*}
\end{example}



\begin{problem}(problem 3)
Compute the average value of $f(x) = \cos(2x) + \sin(2x)$ on the interval $[0,\pi/2]$.
\begin{hint}
An antiderivative of $\cos(2x) + \sin(2x)$ is  $\tfrac12\sin(2x) - \tfrac12\cos(2x)$
\end{hint}
\[f_{ave} = \answer{2/\pi}.\]
\end{problem} 



\begin{example}[example 4]
Compute the average value of $f(x) = e^{2x}$ on the interval $[-1,1]$.\\
The average value is given by 
\begin{align*}
f_{ave} &= \frac{1}{b-a}\int_a^b f(x) \ dx \\
 & = \frac12\int_{-1}^1 e^{2x} \ dx \\
 & = \frac14 e^{2x} \Bigg|_{-1}^1 \\
  & = \frac{e^2 - e^{-2}}{4}
\end{align*}
\end{example}



\begin{problem}(problem 4)
Compute the average value of $f(x) = e^{3x}$ on the interval $[-1,1]$.
\[f_{ave} = \answer{\frac{e^{3} - e^{-3}}{6}}.\]
\end{problem} 

\end{document}
